\documentclass[12pt,onesided,fullpage]{amsart}
\usepackage{amsfonts}
\usepackage{amsmath}
\usepackage{amssymb}
\usepackage{array,amsmath,graphicx,psfrag,amssymb,subfigure,tabularx,booktabs}
\usepackage{mparhack}
\usepackage{setspace}
\usepackage{natbib}
\usepackage{multicol}
\usepackage{dcolumn}
\usepackage{hyperref}
\usepackage{multirow}
\usepackage{color}
\usepackage[top=3cm, bottom=3cm, left=2.3cm, right=2.3cm]{geometry} 
%\pdfpagewidth=8.5in % for pdflatex
%\pdfpageheight=11in % for pdflatex
\begin{document}

\title{Revision Memo For Journal of Politics}
Dear Professor Gailmaird,
Blablabla, what we do is in \textcolor{blue}
Best,
Cassy Dorff, Max Gallop, and Shahryar Minhas

\subsection{Editor's Comments}
Dear Professor Dorff,

I have received the reviews of your manuscript, appended at the end of this email, and read the paper myself. I am happy to report that, based on these evaluations, I am offering you the opportunity to revise and resubmit your manuscript for further consideration at the JOP. 

The reviews, as you can see, are all fairly positive. They believe that your approach brings fresh insights to the literature on civil war, and that there is an important contribution here. At the same time, they are also all fairly critical about the paper's clarity in presenting its central aims and defending the analytical strategy for addressing those claims. I found the reports convincing, both in their praise for the paper and in their critiques. While there is a fair amount of work to do in light of these reports, I think it is eminently doable. 

Though this piece came to me as a methods paper, I agree with R. 2 that the best approach is probably to focus foremost on the substantive motivation, lay out your original theoretical perspective, show how that determines the choice of methodology, and explain what you found. That dovetails with R. 3's good suggestions about approaching the presentation, and with R. 1's critiques about the clarity of the central aim of the paper. The key questions are boiled down succinctly by R. 3 who asks that you clarify, "what are the consequences of these limitations [that you identify in previous approaches]? What don't we understand that, at the end of this article, we will understand better?"

\textcolor{blue}{We clarify the purpose of the paper, principally in }


In addition, though you certainly sought to avoid wading into the latent space "vs." (T)ERGM "debate," I agree with R. 1 that more discussion of your preference for one over the other is called for. To me the most straightforward (if labor intensive), and most transparent way to handle this is to consider TERGM as a robustness check in an appendix. Alternatively, if you have a strong reason not to do so, you should better justify this with citations to the literature and/or theoretical discussion for your case. In essence, choosing only one approach and downplaying the other is inherently to take sides in the 'debate,' but a paper that hands primarily on a substantive point is not the right place to do so. 

\textcolor{blue}{In the technical appendix, we reestimate these models using a TERGM. The performance is BLAH and the ...}

The referees offer a number of other thoughtful critiques and suggestions, all of which you should address in your revision memo detailing whether and where in the revision you followed them, and explaining why you did not if you choose not to do so. 

When I receive the revised manuscript I will send it back to these three reviewers. I will send it to a wholly new reviewer only if one of the original reviewers is unavailable to read the revision.

Additional instructions on the JOP R&R process are below the signature line of this email. Note that while the invitation to revise and resubmit the manuscript does not constitute a commitment to publish, I am aware of the considerable time and effort required to undertake revisions. I would not invite you to do so if I was not confident that you can satisfy the reviewers' concerns and gain a favorable decision. 

I am convinced that this is interesting and important work with the potential to make a significant contribution to the field. I look forward to reading your revised manuscript.

Sincerely,
Sean Gailmard
Field editor, Formal theory and methodology 
JOP


Comments from the Reviewers:

Reviewer #1: I greatly enjoyed reading this paper. It offers a number of important contributions. First, it provides an empirical modeling framework for research on the internal dynamics of civil conflict that account for the complex interdependent structure underlying the process under study. Second, it extends existing latent variable models for networks to account for exogenous over-time variation in the set of active vertices. Third, it demonstrates that their network model out-performs a simple GLM in predicting conflict. Fourth, the results offer important insights regarding the drivers of civil conflict, and illustrate how the proposed modeling framework can be used to understand the effects of individual actors. This will be an important foundational paper for empirical analyses of violent (and non-violent) interactions within civil conflict. All this said, I think this paper requires some revisions before it is ready for publication. I detail three directions for
revisions below.

First, the authors offer a lot in this paper, but it is difficult to evaluate their contributions in the context of their research objective. The research objective of this paper should be clarified. I can see four different aims that are addressed at least partially in this paper. The first is to develop the best performing classifier to predict battles in a civil conflict. The second is to understand whether a latent variable network model performs better than GLM in predicting battles, with the intent of interpreting the performance gain as the benefit of accounting for network dependencies. The third is to develop an explanatory model of battles from which we can learn about the factors that serve as indicators of civil conflict. The first is to propose and illustrate the application of a latent variable network model that can incorporate exogenous variation in the vertex set. Once the authors clarify their objective, it would be useful to provide additional analysis
that builds toward their aim. For example, if the goal is to develop an effective classifier, it would be valuable to see comparisons to random forests, neural networks, and other learning approaches at are designed to optimize predictive performance. I realize that these classifiers may not seem as interpretable, but recent developments render interpretation much more feasible (see, e.g., Jones and Lupu (forthcoming in AJPS)).

\textcolor{blue}{In the appendix, we compare our models performance to a number of binary classifiers and find...}


Second, the empirical modeling presented in this paper is generally very carefully executed. One aspect of the model specification that I would recommend exploring is the degree to which there is temporal heterogeneity in the parameter values. I am not an expert in civil conflict, but it seems likely that the factors governing battle instances will shift over time throughout the conflict. I think the authors have plenty of data to estimate parameters separately for each time period. It would be useful to see plots of the parameter credible intervals over time to understand the degree of temporal heterogeneity in the relationships between features and battle instances. This would be especially useful if the objective of the paper is focused on learning about the factors that drive battles in civil conflict.

\textcolor{blue}{We run a separate model where we estimate separate AME models for each time period. For obvious reasons, we exclude the variable for whether a time period is after the entrance of Boko Haram. The results are included in the appendix.} 

The third consideration regards the modeling framework used. The authors, without referencing any literature, cite the need for the ERGM (and presumably the temporal ERGM (TERGM)) to treat the set of actors as constant over time. This is the only reason they present for not using the TERGM as their modeling approach. It is not true that TERGM requires a constant actor set. Actor entry and exit in the TERGM is exogenous, as it is in the model proposed by the authors. The authors note, "we adjust the estimation procedure of the AME model such that actor observations only contribute to the likelihood of the model once they enter into the network." This is precisely how this is handled with TERGM. Indeed, most of the political science applications of TERGM with which I am aware do not subset the nodes to get a constant active node set. The Journal of Statistical Software article on btergm (https://www.jstatsoft.org/article/view/v083i06/v83i06.pdf) explains this---see Section 5.
This said, I recommend that the authors add the TERGM as a comparison model, or provide additional justification for not comparing their model to TERGM. For example, perhaps after the goals of the paper are clarified, it would be clear that there would be no need to compare their model to yet another network model that can account for dependence. TERGM represents an alternative framework for accounting for the same sorts of dependence that are modeled with the latent variable models, and can also incorporate variation in the vertex set. TERGM can be formulated to account for first order dependencies via either in-star and out-star terms and/or lagged covariates that measure the number of ties sent and/or received by the sender and/or recipient in the past. Higher order dependencies can be modeled using edgewise and dyad-wise shared partner terms  and/or delayed shared partner terms that measure the completion of closed and open triads over time. At the very least, it would
be great to see a robustness check in which the authors evaluate whether the effects of covariates are robust to modeling with TERGM.

\textcolor{blue}{As stated above, we compare our model to a TERGM in the appendix.}

I have two more minor points---I would like to see a discussion of setting the priors used for inference with this model. In particular, how does the specification of the priors relate to the predictive performance of the latent variable model? If the priors make any sort of difference, how should the priors be set to assure that the model performs as well as possible while also assuring that the model is not trained on the test data. 

\textcolor{blue}{We discuss the priors on page ...}

Second, for the predictive experiment, it would be great to see if the comparison of the AME and GLM is sensitive to k---the number of groups into which the data is split. The latent variable structure in the AME renders the observations within a time period dependent upon each other. This could be seen in the posterior predictive distribution if it were possible to integrate the joint likelihood of the observed and unobserved data over the priors. As such, which observations are included in the training set provides some information about the
observations excluded from the training set. The independent information contained in the held out data should increase as k decreases. I couldn't find the value of k in the manuscript, but it would be great to see if the predictive performance of AME relative to GLM varied between, say k=5 and k=20.

\textcolor{blue}{We insert a mention of the value of $k$ used in our analysis on page X. We also, in the appendix, compare the predictive value for values of $k$ ranging from 5 to 20.}

Reviewer #2: The manuscript explores the effect of actor interdependence on civil war dynamics. The question is new and interesting, and the authors have some intriguing answers. I like the substantive story and think that this paper has a great potential to make a significant contribution. The biggest problem is the current presentation.  The good news is that, although it would require a lot of editing and re-writing, the presentational issues are ultimately superficial and can be fixed.

\textcolor{blue}{We agree that our prose is unclear, and we have made an effort to clarify the aims of the paper and generally improve the writing.}

What I find interesting and publication-worthy about this paper is the substantive insights about actors, their effects of one another, and the entry of new conflict participants. This is the part that constitutes the core of the contribution, not the methodological approach. I would support publication of this manuscript as a substantive piece that is driven by an empirical puzzle and analyzed using an appropriate method.  The manuscript, as it is currently written, sells itself as methodological fix rather than a substantive advance on our knowledge of civil war dynamics. I think this is a poor framing/presentational choice and sells the contribution short. It also seems backwards.  I would start with a substantive puzzle, build a theoretical model of interaction among actors of different types, derive predictions, and only then get to the method and testing. As it is, the actual substantive story does not start until p. 8--this is where I would start the paper.

\textcolor{blue}{We think this is very useful in terms of the framing of the paper, and we make efforts to do so.}

I especially like the part about new participant entry. Why would we expect that an entry by a new participant would affect the conflict dynamics?  What types of participants are likely to enter an ongoing conflict? What possible relationships may a new participant have to the current participants?  Why would an entry of a new participant increase outbidding?  I think the authors have some interesting ideas, but these need to be developed much further. I would like to be presented with a well-grounded theoretical story related to the entry of a new participant in an ongoing conflict. 

\textcolor{blue}{We discuss the reasons different groups might destabilize conflict dynamics on page XXX.}

Instead, the current draft spends the first 7 pages trying to convince the reader that the study of civil war is important and that actors in a civil war may affect one another's actions--both of these ideas are non-controversial. Then there is a lengthy discussion of network terms and methods, which I would move to the research design section.  I would also move the history of Nigerian conflict to the appendix, is distracts from the flow.  As it is, various parts of the paper read as disconnected from one another, almost as if they are copied/pasted from multiple different papers.

\textcolor{blue}{We maybe move this section to the appendix.}

Reviewer #3: This manuscript seeks to model intrastate conflict through network analysis and thus to more directly examine interdependencies among actors, as well as the potential for changes in actors over time. It seeks to build on a growing literature on multiparty conflicts and to move beyond dyadic research designs (which were themselves an advance over country-level monadic studies). There is a lot of good material here, and I really like the intuition behind this manuscript. I see it as having great potential. However, as it stands, I do not find that the manuscript is presented or argued in a way that makes a clear contribution to knowledge. For the manuscript to be publishable, it would need substantial revision.


It is not entirely clear at any point in this manuscript what problem the manuscript is trying to address. The author(s) are correct that civil war scholars have not properly considered the interdependencies among actors in multiparty conflicts and that dyadic research designs are limited. However, what are the consequences of these limitations? What don't we understand that, at the end of this article, we will understand better?

\textcolor{blue}{We agree that the paper was unclear, we have tried to clarify the problems this manuscript was trying to address.}

The lack of focus continues through the manuscript. The author(s) state on page 3 "This study investigates how relationships between armed actors evolve over time, and
how the structure of relationships in one time period influences the occurrence of violence
between armed actors in a future period. To do this, we consider how interdependence
across actors, changes in actor composition over time, and actor-level attributes predict
violence in the network." There are several potential questions here, such as whether different armed groups fight (or cooperate), how the presence of additional actors affects the level of violence between rebel groups and the state, the overall level of violence in the conflict, etc. The lack of specificity comes through in the theory, where there is a discussion of key concepts in network theory like actor centrality and reciprocity, but it is unclear how exactly these concepts relate to the question in the manuscript or to existing work. It is also present in the analysis, where the DV is the occurrence of battles between armed actors, but it is not clear why that is appropriate given the theory.

\textcolor{blue}{We appreciate the feedback on the need for clarity in a number of areas. We especially focus on the question XXXX. We attempt to clarify how network dynamics relate to existing work on page XXX, and our choice of a DV on XXX.}

All of this is fixable, but it requires re-framing the introduction, better situating the argument in the literature, and re-writing the theory. Why not start with the question?why do some armed groups fight each other and others do not? There is clear literature on this question, much (although not all) of which is cited here. Describe that literature and explain why it is insufficient to understand this phenomenon. Then, the manuscript needs a clear, focused theoretical argument drawing on network theory to explain this. The theory should start from first principles?what are armed actors trying to achieve through violence and how do the components here (centrality, reciprocity, the introduction of new "aggressive" actors, etc.) fit into this story.

\textcolor{blue}{We reframe the introduction, and focus on the behavior of armed actors, in particular blablabla.}

With this set-up, the analysis would then be much more closely tied to the theory. The manuscript also needs much more discussion of how the results of the analysis extend our knowledge. What do they tell us about existing work? What do they contribute to what we already know?

\textcolor{blue}{We discuss how this helps inform our understanding of existing work on...}

This is my main comment, that the manuscript needs to make a much clearer case for how introducing the network dynamics described here contribute to knowledge on the question being asked. I have a few, smaller, comments:

*       To identify groups, the author(s) use ACLED. I understand why, but the downside of doing this is that they lose the ability to include group level characteristics present in, for example, the Non-State Actor data, the EPR project, etc. The authors need to make a clearer case for the contribution here given that they are giving up substantial access to information about these actors.

\textcolor{blue}{We make the case for our use of the ACLED data on page XXX.}

*       The manuscript needs to be much clearer on what the system is. Here, it includes everyone fighting in Nigeria. But is that always appropriate? Is the country the box that bounds networks? Why? What about separate, peripheral, conflicts in places like India, should they be modelled together? What about linked conflicts in places like Central Africa? Should they be modelled separately? These considerations are important.

\textcolor{blue}{We discuss how the modeling framework would differ given transnational or subnational conflicts on page XXX.}

*       The manuscript should, at the least, cite Bapat and Bond (2011), which is one of the earliest articles to examine cooperation/conflict between rebel groups in civil war.

\textcolor{blue}{We discuss Bapat and Bond's paper on page XXX.}

\end{document}\byei
